\documentclass[10pt,a4paper]{article}
\usepackage[T1]{fontenc}
\usepackage{lmodern}
\usepackage[margin=18mm]{geometry}
\usepackage{tabularx,booktabs,array}
\setlength{\parindent}{0pt}
\setlength{\parskip}{4pt}
\pagestyle{empty}

\begin{document}

{\small VU MIF \enspace|\enspace Programų sistemų inžinerija\par}
{\LARGE\bfseries Sistemos projekto rengimo planas\par}

\textbf{Tikslas.} Parengti sistemos projektą, pagal kurį programuotojai galėtų
įgyvendinti sistemą be papildomų klausimų apie jos architektūrą, sąsajas,
duomenis ar elgseną. Projektą grindžiame ankstesniais reikalavimų ir analizės
dokumentais, išlaikydami jų identifikatorius ir terminus.

\textbf{Metodika.} Dokumentą struktūruojame pagal 4+1 architektūros vaizdų modelį.
Loginis vaizdas aprašo domeną ir taisykles, kūrimo vaizdas -- modulius ir jų
priklausomybes, procesų vaizdas -- sistemos elgseną vykdymo metu, fizinis
vaizdas -- komponentų diegimą. Panaudojimo atvejai (PA) sudaro „+1“ ir susieja
vaizdus. Modulių ir objektų atsakomybėms apibrėžti taikome objektinį
projektavimą (OOD). Sudėtingesniems PA galime taikyti ICONIX analizės ir
sekų diagramas, susiejančias ekranus, valdymo logiką ir domeno objektus.

\subsection*{Darbų paskirstymas}

\renewcommand{\arraystretch}{1.18}
\begin{tabularx}{\linewidth}{@{}r l >{\raggedright\arraybackslash}X@{}}
\toprule
\textbf{Nr.} & \textbf{Atsakingas} & \textbf{Rezultatas} \\
\midrule
1 & Tadas & \textbf{Įvadas, architektūra ir diegimas.}
Metodika, sistemos konteksto, komponentų ir diegimo diagramos.
Sluoksnių ir modulių atsakomybės bei priklausomybės. Pagrįstas vietinio,
debesijos ar hibridinio diegimo pasirinkimas, DB vieta ir vietinio grotuvo
ryšys su kitais komponentais. \\
2 & Marijus & \textbf{Duomenų bazės projektas.}
Fizinė schema, laukų tipai, raktai, apribojimai ir indeksai.
Transakcijos, archyvavimas, šalinimas ir pagrįstas lentelių skaidymo
(\textit{partitioning}) sprendimas. \\
3 & Goda & \textbf{Naudotojo sąsaja.}
Ekranų žemėlapis, „Figma“ eskizai ir susietas pagrindinių srautų prototipas.
Veiksmai pagal roles, laukų validavimas, tuščios ir klaidų būsenos.
Ekranų nuorodos į PA ir susijusias operacijas. \\
4 & Tomas & \textbf{API ir komponentų sąsajos.}
Sąsajos tarp naudotojo sąsajos, serverio ir vietinio grotuvo pagal
pasirinktą architektūrą. Operacijų įvestis, išvestis, DTO, prieigos teisės
ir klaidos. Pasirinkus HTTP API, jos specifikacija pateikiama „OpenAPI“ formatu. \\
5 & Ugnė & \textbf{Verslo logika, būsenos ir atsekamumas.}
Verslo taisyklės, būsenų perėjimai, pagrindinių ir alternatyvių scenarijų
sekos. Bendra atsekamumo matrica ir jos spragų patikra. \\
\bottomrule
\end{tabularx}

\subsection*{Detalumas ir projekto patikra}

Pradžioje komanda pagal ankstesnius dokumentus suderina modulių ribas, diegimą,
roles ir būsenas. Kiekvienas savo dalyje susieja reikalavimus ir PA su
atitinkamais ekranais, operacijomis, moduliais ar duomenimis bei priėmimo
kriterijais. Ugnė sujungia šias nuorodas į bendrą matricą.

Paprastam veiksmui pakanka taisyklės ir sąsajos aprašo. Sudėtingam scenarijui
pridedame sekų arba būsenų diagramą. Pavyzdžiui, rezervuojant paskutinę
vietą turi būti aišku, kas gali rezervuoti, kokie duomenys perduodami,
kur tikrinamos laisvos vietos, kas vyksta gavus du prašymus vienu metu
ir ką sąsaja parodo abiem naudotojams.

\textbf{Baigtumo kriterijus.} Scenarijaus nerengęs komandos narys pagal
dokumentą gali suskaidyti jį į programavimo užduotis ir nustatyti laukiamą
rezultatą sėkmės bei klaidų atvejais. Jei reikia autoriaus paaiškinimo,
trūkstamą sprendimą įrašome į projektą.

\end{document}
